\documentclass[17pt, a0paper, landscape]{tikzposter}

% =============================================================================
% Standard Packages
% =============================================================================
\usepackage{amsmath, amssymb}
\usepackage{booktabs}
\usepackage{graphicx}
\usepackage{microtype}
\usepackage{xcolor}

% Automatically search both root and results/ directories
\graphicspath{{./}{results/}{../results/}}

% Fail-Safe Image Macro: Pure LaTeX boxes (No nested TikZ, detokenized underscores)
\newcommand{\SafeImage}[3][width=0.96\linewidth]{%
  \IfFileExists{#2}{%
    \includegraphics[#1]{#2}%
  }{%
    \IfFileExists{results/#2}{%
      \includegraphics[#1]{results/#2}%
    }{%
      \setlength{\fboxsep}{10pt}%
      \setlength{\fboxrule}{1.5pt}%
      \fcolorbox{ThemeSecondary}{ThemeSecondary!8}{%
        \parbox{0.9\linewidth}{%
          \centering\vspace{1.0em}%
          {\Large\bfseries\color{ThemePrimary} #3}\\[0.6em]%
          {\normalsize\ttfamily\color{ThemeSecondary} [Graphic Placeholder: \detokenize{#2}]}\\[0.4em]%
          {\small\itshape\color{gray} (Upload image to Overleaf to display graphic)}%
          \vspace{1.0em}%
        }%
      }%
    }%
  }%
}

% =============================================================================
% Modern Theme & Refined Nordic Slate & Teal Palette
% =============================================================================
\usetheme{Simple}
\usecolorstyle{Default}

\definecolor{ThemePrimary}{HTML}{0F172A}    % Slate 900 / Deep Midnight
\definecolor{ThemeSecondary}{HTML}{1E293B}  % Slate 800 / Tech Navy
\definecolor{TealAccent}{HTML}{0D9488}      % Modern Deep Teal
\definecolor{CyanAccent}{HTML}{0284C7}      % Electric Sky Cyan
\definecolor{EmeraldGreen}{HTML}{059669}    % Success Green
\definecolor{AmberWarning}{HTML}{D97706}    % Amber
\definecolor{CanvasBg}{HTML}{F8FAFC}        % Clean Off-White Tint
\definecolor{CardBg}{HTML}{FFFFFF}          % Crisp White

\colorlet{backgroundcolor}{CanvasBg}
\colorlet{framecolor}{ThemeSecondary}
\colorlet{titlefgcolor}{white}
\colorlet{titlebgcolor}{ThemePrimary}
\colorlet{blocktitlebgcolor}{ThemeSecondary}
\colorlet{blocktitlefgcolor}{white}
\colorlet{blockbodybgcolor}{CardBg}
\colorlet{blockbodyfgcolor}{ThemePrimary}

% =============================================================================
% Title, Author & Metadata
% =============================================================================
\title{\textbf{Pavement Crack Segmentation with Symmetrical U-Net}}
\author{\textbf{Tackling Severe Class Imbalance (60.5 : 1) via Validation-Tuned Weighted Cross-Entropy}}
\institute{Machine Learning Project $\bullet$ Problem Statement 2 $\bullet$ Benchmark: CrackForest Dataset (CFD) $\bullet$ 1.94M Parameters}

% =============================================================================
\begin{document}
% =============================================================================

\maketitle

% =============================================================================
% FULL-WIDTH TOP BANNER: End-to-End Pipeline SmartArt
% =============================================================================
\block[titleinnersep=5pt,bodyinnersep=5pt]{End-to-End Methodology \& Pipeline Workflow (SmartArt Process Flow)}{
    \begin{center}
        \SafeImage[width=0.98\linewidth]{pipeline_smartart.png}{End-to-End Deep Learning Pipeline (Ingestion -> Architecture -> Loss -> Tuning -> Inference)}
    \end{center}
}

\begin{columns}

% =============================================================================
% COLUMN 1: Motivation, Dataset & Symmetrical U-Net Architecture SmartArt
% =============================================================================
\column{0.333}

\block{1. Problem Motivation \& 60.5 : 1 Imbalance}{
    \textbf{Real-World Challenge:} Continuous automated detection of asphalt cracks is essential for structural civil maintenance and proactive pothole prevention. However, cracks are ultra-thin spatial anomalies occupying only 1--5 pixels in width.

    \vspace{0.3em}
    \textbf{CFD Pixel-Level Imbalance:}
    \begin{itemize}
        \item \textbf{Crack Foreground:} 292,885 pixels (\textbf{1.63\%} of total surface area).
        \item \textbf{Pavement Background:} 17,710,362 pixels (\textbf{98.37\%} of total surface area).
        \item \textbf{Imbalance Ratio:} \textbf{60.47 : 1} background to crack pixels.
    \end{itemize}

    \vspace{0.3em}
    \textbf{Why Standard BCE Fails (The Accuracy Trap):}
    Under standard unweighted BCE, gradients from the overwhelming $>98\%$ background pixels dominate backpropagation. A degenerate network predicting 100\% background achieves \textbf{98.37\% Pixel Accuracy} while detecting 0\% of cracks. Standard BCE produces broken, dotted lines or omits entire fissures (Baseline Recall was only 37.84\%).
}

\block{2. Dataset Curation \& Controlled Splits}{
    \textbf{CrackForest Dataset (CFD):} 118 calibrated color road images ($\sim 480 \times 320$) paired with pixel-level ground truth binary masks.
    
    \vspace{0.2em}
    \begin{itemize}
        \item \textbf{Standardized Resolution:} Resized to \textbf{256 $\times$ 256} (Bilinear for RGB; Nearest-Neighbor for masks to preserve $\{0, 1\}$ binary labels).
        \item \textbf{Deterministic Split (Seed = 42):}
        \begin{itemize}
            \item \textbf{Train Split:} 82 images (69.5\%) -- weight optimization.
            \item \textbf{Validation Split:} 17 images (14.4\%) -- hyperparameter selection of $w$.
            \item \textbf{Test Split:} 19 images (16.1\%) -- held-out final benchmark.
        \end{itemize}
        \item \textbf{Augmentation:} Random horizontal and vertical flips ($p=0.5$).
        \item \textbf{In-Memory Caching:} All 118 images pre-cached in RAM for high throughput.
    \end{itemize}
}

\block{3. Symmetrical U-Net Architecture (SmartArt)}{
    \begin{center}
        \SafeImage[width=0.96\linewidth]{unet_architecture_smartart.png}{Symmetrical U-Net: 4 Encoders, Skip Connections, Bottleneck \& 4 Decoders}
    \end{center}
    
    \vspace{-0.3em}
    \begin{itemize}
        \item \textbf{Encoder (Contracting):} 4 stages [DoubleConv 3$\times$3 $\rightarrow$ BN $\rightarrow$ ReLU $\rightarrow$ MaxPool 2$\times$2]. Channels: $3 \rightarrow 16 \rightarrow 32 \rightarrow 64 \rightarrow 128$.
        \item \textbf{Bottleneck:} 256 feature channels at $16 \times 16$ spatial resolution.
        \item \textbf{Decoder (Expanding):} 4 stages [ConvTranspose2d stride 2 + DoubleConv]. Channels: $256 \rightarrow 128 \rightarrow 64 \rightarrow 32 \rightarrow 16$.
        \item \textbf{Skip Connections (Red Arrows):} Concatenates fine coordinates directly from encoder to decoder, preventing thin 1-pixel cracks from vanishing.
        \item \textbf{Policy:} Adam ($\text{lr} = 1.5 \times 10^{-3}$, decay $= 10^{-5}$), CosineAnnealingLR.
    \end{itemize}
}

% =============================================================================
% COLUMN 2: Loss Formulation, Multi-Metric Tuning & Visual Proof
% =============================================================================
\column{0.334}

\block{4. Loss Formulation \& Gradient Derivation}{
    \textbf{Ordinary Binary Cross-Entropy (Baseline, $w = 1.0$):}
    \begin{equation*}
        \mathcal{L}_{\text{BCE}} = -\frac{1}{N} \sum_{i=1}^N \left[ y_i \log(p_i) + (1 - y_i) \log(1 - p_i) \right]
    \end{equation*}

    \textbf{Proposed Weighted Binary Cross-Entropy (WBCE, $w > 1.0$):}
    \begin{equation*}
        \mathcal{L}_{\text{WBCE}}(w) = -\frac{1}{N} \sum_{i=1}^N \left[ w \cdot y_i \log(p_i) + (1 - y_i) \log(1 - p_i) \right]
    \end{equation*}
    where $y_i \in \{0, 1\}$ and $p_i = \sigma(z_i) = \frac{1}{1 + e^{-z_i}}$.

    \vspace{0.3em}
    \textbf{Gradient Derivation on Falsely Neglected Cracks:}
    \begin{equation*}
        \frac{\partial \mathcal{L}_{\text{WBCE}}}{\partial z_i} = p_i - y_i \left[ 1 + (w - 1)(1 - p_i) \right]
    \end{equation*}
    When a true crack pixel is falsely predicted as pavement ($y_i = 1, p_i \approx 0$), the error gradient is multiplied by \textbf{$w$}, applying heavy corrective backpropagation.

    \vspace{0.2em}
    \textbf{Controlled Ablation:} Setting $w = 1.0$ mathematically recovers baseline BCE, isolating the exact empirical contribution of class reweighting.
}

\block{5. Multi-Metric Tuning: F1, Accuracy, Precision \& Recall vs $w$}{
    \begin{center}
        \SafeImage[width=0.96\linewidth]{validation_weight_tuning.png}{Validation Tuning Curves: F1, Accuracy, Precision, Recall \& IoU}
    \end{center}
    
    \vspace{-0.3em}
    \textbf{Comprehensive Multi-Metric Analysis:} 
    \begin{itemize}
        \item \textbf{Pixel Accuracy (brown):} Remains stably high ($>97.5\%$) across all weights due to background volume.
        \item \textbf{Recall (orange):} Climbs monotonically from $45.28\%$ ($w=1$) to $90.05\%$ ($w=15$).
        \item \textbf{Precision (red):} Gradually trades off as weight increases ($67.23\% \rightarrow 39.13\%$).
        \item \textbf{F1 Score (blue) \& IoU (green):} Peak at \textbf{$w^* = 5.0$} with \textbf{Val F1 = 0.6153} and \textbf{Val IoU = 0.4444}, establishing $w^* = 5.0$ as the optimal trade-off point.
    \end{itemize}
}

\block{6. Qualitative Visual Continuity on Test Images}{
    \begin{center}
        \SafeImage[width=0.96\linewidth]{test_predictions_comparison.png}{Qualitative Predictions: Input vs Ground Truth vs Baseline vs Proposed}
    \end{center}
    
    \vspace{-0.3em}
    \textbf{Visual Finding:} Baseline U-Net (Col 3) yields broken, fragmented fissures. The proposed model with $w^* = 5.0$ (Col 4) completely reconstructs unbroken crack continuity.
}

% =============================================================================
% COLUMN 3: Benchmark Results, Confusion Matrix & Defense
% =============================================================================
\column{0.333}

\block{7. Quantitative Benchmark: F1, Accuracy, Precision \& Recall}{
    \begin{center}
        \SafeImage[width=0.96\linewidth]{metrics_comparison_bar.png}{Test Benchmark: Pixel Acc, Precision, Recall, F1 Score \& IoU}
    \end{center}
    
    \vspace{-0.2em}
    \begin{center}
    \small
    \begin{tabular}{lcccc}
        \toprule
        \textbf{Metric} & \textbf{Baseline ($w=1$)} & \textbf{Proposed ($w=5$)} & \textbf{$\Delta$ Gain} & \textbf{Relative} \\
        \midrule
        \textbf{Pixel Accuracy} & 98.77\% & 98.37\% & -0.40\% & -- \\
        \textbf{Precision} & 67.97\% & 48.17\% & -0.1980 & -29.1\% \\
        \textbf{Crack Recall} & 37.84\% & \textbf{78.09\%} & \textbf{+0.4025} & \textbf{+106.4\%} \\
        \textbf{F1 Score (Dice)} & 0.4862 & \textbf{0.5958} & \textbf{+0.1097} & \textbf{+22.6\%} \\
        \textbf{IoU (Jaccard)} & 0.3211 & \textbf{0.4243} & \textbf{+0.1032} & \textbf{+32.1\%} \\
        Mean Class Acc & 68.78\% & 88.39\% & +19.61\% & +28.5\% \\
        \bottomrule
    \end{tabular}
    \end{center}

    \vspace{0.2em}
    \textbf{The Accuracy Paradox:} High pixel accuracy ($98.77\%$) in the baseline is misleading because it missed $62.16\%$ of cracks. The proposed model doubled crack recall ($78.09\%$), delivering massive real-world inspection gains (+22.6\% F1, +32.1\% IoU).
}

\block{8. Pixel-Level Confusion Matrix Analysis}{
    \begin{center}
        \SafeImage[width=0.96\linewidth]{confusion_matrix.png}{Confusion Matrix: 1,245,184 Held-Out Test Pixels}
    \end{center}
    
    \vspace{-0.3em}
    \begin{itemize}
        \item \textbf{64.8\% Reduction in Missed Cracks:} False Negatives dropped from 11,888 pixels (Baseline) down to only 4,190 pixels (Proposed).
        \item \textbf{Controlled Background FP:} Pavement false positives remained at only 1.31\%, preserving high overall pixel reliability ($98.37\%$).
    \end{itemize}
}

\block{9. Technical Viva Q\&A \& Key Takeaways}{
    \begin{itemize}
        \item \textbf{Q: Why did ordinary BCE produce fragmented cracks?} \\
        \textbf{A:} Because 98.4\% of pixels are pavement, backprop gradients pushed predictions towards 0 to minimize aggregate loss, ignoring thin foreground signals.
        \item \textbf{Q: Why not set $w = 60.5$ (inverse frequency)?} \\
        \textbf{A:} Setting $w = 60$ excessively penalizes false negatives, causing road texture and shadows to be flagged as cracks. Validation tuning proved $w^* = 5.0$ is the optimal sweet spot.
        \item \textbf{Q: Why are skip connections indispensable?} \\
        \textbf{A:} MaxPool downsamples 1-pixel cracks to zero signal. Skip connections re-inject high-frequency spatial coordinate details directly to the decoder.
    \end{itemize}
}

\end{columns}

\end{document}
